\documentclass[12pt]{article}
\usepackage[utf8]{inputenc}
\usepackage[T1]{fontenc}
\usepackage{hyperref}
\title{Yulu Language Resources for K-12 Curriculum Development}
\author{Compiled by Claude Sonnet 5.5 (claude-sonnet-5-5) for Andr\'as Kornai}
\date{}
\begin{document}
\maketitle

\section{Introduction}
Yulu (ISO 639-3 yul; Glottocode yulu1243) is a Central Sudanic language of the Central African Republic and neighbouring Sudan/South Sudan. Search results give estimates of 7,000 to 13,000 speakers and describe it as vulnerable or endangered. No evidence was found that it is taught in schools. Very little published material exists, and nothing found is suitable for classroom use.

\section{Primary Educational Resources}
No primer, textbook or school programme for Yulu turned up in the searches. A curriculum would have to start from scratch, using the descriptive sources below.

\section{Linguistic and Reference Materials}
\begin{itemize}
\item Wikipedia overview of Yulu: \url{https://en.wikipedia.org/wiki/Yulu_language}. It mentions the early work of Stefano Santandrea.
\item Endangered Languages Project entry for Yulu, which lists linguistic publications including Pascal Boyeldieu's work on compound verbs and a 1976 St Andrews thesis on Yulu phonology by James Gabjanda: \url{https://endangeredlanguages.com/lang/2307}
\item Pascal Boyeldieu's chapter on similative and equative constructions, which includes Yulu: \url{https://benjamins.com/catalog/tsl.117.10boy}
\item Yulu numeral list from the Max Planck EVA numerals database: \url{https://lingweb.eva.mpg.de/channumerals/Yulu.htm}
\item Library catalogue record for a 1970 book of brief grammar outlines of Yulu and Kara, with a small comparative vocabulary: \url{https://library.aiu.ac.ke/bib/64160}
\end{itemize}
No Yulu dictionary was confirmed.

\section{Religious and Community Resources}
None found.

\section{Digital and Audio Resources}
No audio, video, corpus or software resources found.

\section{Curriculum Development Context}
Yulu is documented only in a few linguistic descriptions. Any K-12 effort would need new documentation work with speaker communities, an agreed orthography, and teaching materials written from scratch. Contact with linguists who have published on Yulu would be a sensible first step.

\end{document}